\section{Nonsmooth scalar curvature and the two level currents}
\label{sec:r11-nonsmooth-scalar-curvature}

Keep the original hypotheses and notation of
Section~\ref{sec:r11-level-disintegration}.  Thus $Q$ is any
original symmetric middle-positive separately closed current of
degree $(p,p)$, $p<n$, and $f$ is a locally bounded difference
of classical psh functions.  All identities are identities of
original coefficient-Radon currents on ordinary open sets, tested
against every original compact smooth complementary form.  The
positive level-current theorem of that section is the only new
input to the present extension.

\begin{definition}[Lower and upper slope conventions]
Write
\[
 \begin{split}
 \mathsf L_f^{-}(t)
    &=D(\max(f,t)Q)-\mathbf1_{\{f>t\}}D(fQ),\\
 \mathsf L_f^{+}(t)
    &=D(\max(f,t)Q)-\mathbf1_{\{f\geq t\}}D(fQ).
 \end{split}
 \tag{SC1}
\]
The superscripts indicate the lower and upper values assigned to
the derivative of the scalar hinge at its corner; in particular
$\mathsf L_f^{-}(t)=\mathsf L_f^Q(t)$ from (LC3).
The background $Q$ is suppressed in this section's notation only.
\end{definition}

\begin{proposition}[Two positive level currents and their jump]
\label{prop:r11-two-level-currents}
Both currents in (SC1) are middle-positive and carried by
$\{f=t\}$.  They satisfy
\[
 \mathsf L_f^{+}(t)=\mathsf L_{-f}^{-}(-t),\qquad
 \mathsf L_f^{-}(t)-\mathsf L_f^{+}(t)
                      =\mathbf1_{\{f=t\}}D(fQ).
 \tag{SC2}
\]
For every original compact smooth test $\eta$, the scalar function
$t\mapsto\langle\mathsf L_f^{-}(t),\eta\rangle$ is right
continuous and has bounded variation on each compact level interval.
Its left limit at every $t$ is
$\langle\mathsf L_f^{+}(t),\eta\rangle$.  On any compact real
interval of levels, the local positive coefficient masses of both
families are uniformly bounded.  Both admit the energy coarea
identity (LC7), with the same energy current $\Gamma_Q(f)$.
\end{proposition}

\noindent\textbf{Status: \texttt{NEW\_PROVED}}.

\begin{proof}
The elementary identity
$\max(-f,-t)=\max(f,t)-f-t$ and $D(tQ)=0$ give the first
formula in (SC2); Theorem~\ref{thm:r11-positive-level-current}
applied to $-f$ proves positivity and the carrier of
$\mathsf L_f^{+}(t)$.  Subtracting (SC1) proves the second formula.

For a fixed compact test, the function
$t\mapsto\langle D(\max(f,t)Q),\eta\rangle$ is Lipschitz:
the potentials differ by at most $|t-s|$, and the fixed signed
coefficient measure $\mu_{Q,D\eta}$ has finite variation on the
test support.  Continuity from below and above for the fixed
signed coefficient measures of $D(fQ)$ gives
\[
 \begin{split}
 \lim_{s\downarrow t}\mathbf1_{\{f>s\}}D(fQ)
       &=\mathbf1_{\{f>t\}}D(fQ),\\
 \lim_{s\uparrow t}\mathbf1_{\{f>s\}}D(fQ)
       &=\mathbf1_{\{f\geq t\}}D(fQ),
 \end{split}
\]
in local coefficient total variation.  This proves the stated
one-sided continuity after combining the two terms of (SC1).
The first term is Lipschitz in the level, and the second is a
tail of a fixed finite signed measure after pushforward by $f$;
their sum therefore has bounded variation on compact intervals.

For the uniform mass assertion, fix a compact original moment
stage and a positive trace control $\chi\beta^{n-p-1}$, where
$\beta$ is the sum of the elementary positive moment squares.
The closed current $Q'=\beta^{n-p-1}\wedge Q$ has degree
$(n-1,n-1)$ and
$\mathsf L_f^{\pm,Q'}(t)=\beta^{n-p-1}\wedge
\mathsf L_f^{\pm,Q}(t)$ by the bounded smooth-slot identities.
The bounded-variation density bound (LC12), applied to $f$ and
$-f$, bounds the two positive trace masses uniformly in $t$.
Finite original coefficient domination proves the assertion for
every compact stage, including any actual moment refinement.

Finally the difference in (SC2), integrated in $t$ against a
bounded Borel scalar function, is zero: for each point $x$ the
indicator of the singleton $\{f(x)\}$ has Lebesgue integral
zero, and Fubini applies to each fixed finite coefficient
variation of $D(fQ)$.  Therefore (LC7) for the minus convention
also proves it for the plus convention.
\end{proof}

\begin{theorem}[An exact scalar curvature formula with a real measure]
\label{thm:r11-stieltjes-scalar-curvature}
Let $J\subset\mathbf R$ be an open interval, let $I\Subset J$
contain the range of $f$ near the compact test support, and let
$F$ be a finite convex real function on $J$.  Let $\nu_F=F''$
be its positive distributional second-derivative measure, and
write $F'_-,F'_+$ for its finite one-sided derivatives on $I$.
Then $F(f)$ is locally a bounded difference of classical psh
functions, and
\[
 \boxed{\quad
 D(F(f)Q)
  =F'_-(f)D(fQ)+\int_J\mathsf L_f^{-}(t)\,d\nu_F(t)
  =F'_+(f)D(fQ)+\int_J\mathsf L_f^{+}(t)\,d\nu_F(t).
 \quad}
 \tag{SC3}
\]
The integrals are weak integrals of coefficient-Radon currents;
they are well defined locally, middle-positive, and carried by
the actual inverse image of a carrier of $\nu_F$.
Only levels in a compact interval near the local range contribute.

The same formula holds for a real difference of two finite convex
functions, with its signed second-derivative measure and the
corresponding signed integral.  All multipliers in (SC3) use the
actual given Borel values.  The construction commutes with
ordinary restriction and every actual moment refinement.
\end{theorem}

\noindent\textbf{Status: \texttt{NEW\_PROVED}}.

\begin{proof}
Write $f=a-b$ on the chosen witness stage.  On a compact interval
strictly larger than $I$ and contained in $J$, the convex function
$F$ is Lipschitz.  Extend its restriction affinely past the two
endpoints using its interior one-sided slopes; this gives a global
convex Lipschitz extension $\widehat F$ agreeing with $F$ on $I$.
Choose $L$ greater than the absolute values of its slopes.  On the
compact range box of $(a,b)$ the function
\[
                  (s,t)\longmapsto\widehat F(s-t)+L(s+t)
 \tag{SC4}
\]
is convex and nondecreasing in each variable.  The frozen
convex nondecreasing composition rule makes (SC4), composed
with $(a,b)$, a bounded classical psh function.  Subtracting
$L(a+b)$ gives a bounded-difference presentation of $F(f)$.

Choose $c<\inf I\leq\sup I<d$ with $[c,d]\Subset J$.
The elementary real convex representation, on $I$, is
\[
 F(s)=A+Bs+\int_{(c,d)}(s-t)_+\,d\nu_F(t),\qquad
 F'_-(s)=B+\nu_F((c,s)).
 \tag{SC5}
\]
The constants $A,B$ absorb the affine contribution of levels
outside $(c,d)$.  The measure in (SC5) is finite.  The
potentials $(f-t)_+$ are uniformly bounded on the compact
stage, so Fubini against the fixed coefficient measures of
$Q$ and distributional differentiation give
\[
 D(F(f)Q)=B D(fQ)
           +\int_{(c,d)}D((f-t)_+Q)\,d\nu_F(t).
 \tag{SC6}
\]
The positive level masses are uniformly bounded by
Proposition~\ref{prop:r11-two-level-currents}; their compact
test pairings are Borel.  They consequently have a well-defined
weak integral with local coefficient-Radon control.  Insert
$D((f-t)_+Q)=\mathbf1_{\{f>t\}}D(fQ)+\mathsf L_f^-(t)$
in (SC6), and apply Fubini to the fixed signed coefficient
measures in its first term.  Its multiplier is exactly
$B+\nu_F((c,f))=F'_-(f)$.  This proves the first equality
in (SC3).

Replacing the strict indicator by the weak indicator gives
$B+\nu_F((c,f])=F'_+(f)$ and the plus level current, proving
the second equality.  This also checks explicitly the jump
contribution at a real atom of $\nu_F$: changing slope
conventions changes the curvature multiplier by
$\nu_F(\{f\})D(fQ)$ and changes the level integral by its
negative, through (SC2).

For any Borel carrier $A$ of $\nu_F$, the positive scalar
kernel construction in (LC9), applied to the finite measure
$\nu_F$ in place of Lebesgue measure for measurability only,
makes the weak integral a positive measure kernel.  Each level
current is carried by $\{f=t\}$, so its integral is carried
by $f^{-1}(A)$.  No assertion that this inverse image is
closed is required.  Finite polarization gives the same
carrier assertion for every original coefficient.
Linearity treats a difference of convex scalar functions.
Intrinsic definitions and the already established refinement
invariance of each level current prove the final assertions.
\end{proof}

\begin{corollary}[Original Kato inequalities and the exact defect]
\label{cor:r11-original-kato}
For $f$ as above, put $f_+=\max(f,0)$ and use
$\operatorname{sgn}(0)=0$.  Then
\[
 \begin{split}
 D(f_+Q)-\mathbf1_{\{f>0\}}D(fQ)&=\mathsf L_f^-(0)\geq0,\\
 D(|f|Q)-\operatorname{sgn}(f)D(fQ)
                 &=\mathsf L_f^-(0)+\mathsf L_f^+(0)\geq0.
 \end{split}
 \tag{SC7}
\]
Both defects are positive coefficient-Radon currents carried by
the actual zero set of $f$.  The same identities hold after
every original positive coefficient trace and against every
original complementary test.  Below top degree the defects
need not be closed.
\end{corollary}

\noindent\textbf{Status: \texttt{DIRECT\_COROLLARY}} of the new
two-level theorem.

\begin{proof}
The first identity is (SC1).  Apply it to $-f$ and add, using
$|f|=f_++(-f)_+$ and
$\mathsf L_{-f}^-(0)=\mathsf L_f^+(0)$ from (SC2).
The two strict indicators give $\operatorname{sgn}(f)$ with
the printed value at zero.  Positivity and the carriers were
proved before the identities were added.  Nonclosedness is
allowed by the explicit differential defects in (LC5), and
Proposition~\ref{prop:r11-nonclosed-level-example} prevents a
universal closedness inference.
\end{proof}

\paragraph{What this calculus resolves and what it leaves open.}
Formula (SC3) integrates a real curvature measure against actual
level currents of a bounded psh difference.  It therefore gives
a precise replacement for formal multiplication by a real Dirac
mass in a nonsmooth scalar chain rule.  Its hypotheses do not
define an arbitrary pullback of a cycle through a moment map,
and no relation quotient for analytic delta presentations follows.
Moreover the level currents retain the defects (LC5).  Treating
them as separately closed backgrounds in a second BT iteration
is a \texttt{FAILED\_ATTEMPT}; (LC13)--(LC15) give its concrete
failure already on the original two-torus.

